\begin{revised}
\subsection{Additional Controlled Experiments}
\label{sec:revision-experiments}

The following experiments address retraining variability, the unused log loss, greedy commitments, structural transfer, real log coverage, and comparable timing. They supplement the reference results above. Inputs and evaluation budgets were fixed before evaluating the new checkpoints. No test result controls checkpoint selection or a subsequent training run. Computation uses an Intel Xeon Platinum 8468 CPU, Python 3.12.13, PyTorch 2.10.0, and PM4Py 2.7.23.6, with one intra operation and one inter operation thread. This differs from the original CPU, ten-thread configuration, and PM4Py 2.7.23.2. Quality evaluations allow five seconds per candidate and 30 seconds for exact search, with a 31-second external cutoff for the latter. Prefix depth remains eight; completion depth is 32 on the original test set and 64 on the added inputs. Every reported candidate cost is independently replayed.

\subsubsection{Retraining, Log-Loss Ablation, and Candidate Repair}

Two new full-objective runs use seeds 17 and 29, following Section~\ref{sec:optimization}. Table~\ref{tab:revision-seeds} reports their selected epochs and results alongside seed 13 and the optional repair decoder. The unguided control needs no training. The same 512 test rows are used in every comparison. Variation across the three full-objective checkpoints summarizes these particular runs; the small seed count and changed execution environment limit population-level conclusions.

% Generated by scripts/render_revision_tables.py.
\begin{table*}[!t]
\begin{revised}
\caption{\begin{revised}Training seed variation and optional repair on the 512-row synthetic test set. Every condition returns 512 legal candidates; all 512 exact references are available. Epoch denotes the validation-selected checkpoint.\end{revised}}
\label{tab:revision-seeds}
\centering
\renewcommand{\arraystretch}{1.15}
\begin{tabularx}{\linewidth}{@{}Yrrrrr@{}}
\toprule
Decoder & Seed & Epoch & \shortstack{Optimal\\count} & \shortstack{Optimal\\(\%)} & \shortstack{Mean\\gap}\\
\midrule
Greedy, unguided & n/a & n/a & 445 & 86.91 & 0.258\\
Greedy, guided & 13 & 49 & 466 & 91.02 & 0.180\\
Greedy, guided & 17 & 27 & 463 & 90.43 & 0.188\\
Greedy, guided & 29 & 48 & 465 & 90.82 & 0.184\\
Repair, unguided & n/a & n/a & 479 & 93.55 & 0.111\\
Repair, guided & 13 & 49 & 498 & 97.27 & 0.039\\
\bottomrule
\end{tabularx}
\end{revised}
\end{table*}


The full-objective checkpoints produce 466, 463, and 465 optimal candidates for seeds 13, 17, and 29. Their mean optimality is $90.76\%$ with sample standard deviation $.30$ percentage points; mean gap is $.184$ with sample standard deviation $.004$. All three improve on the 445 optimal unguided candidates. The gains are 4.10, 3.52, and 3.91 points. Using the same stratified paired family bootstrap as Section~\ref{sec:protocol}, their respective 95\% intervals are $[2.54,5.66]$, $[1.95,5.08]$, and $[2.54,5.27]$ points. These intervals describe generator sampling for each fixed fitted model, while the standard deviation describes the three observed training runs. They should not be interpreted as interchangeable uncertainty measures.

For each new seed, a matched run sets only the log-loss coefficient to zero. Architecture, remaining losses, data, seed, and optimizer settings are unchanged. Each run retains its own validation selection, so the intervention includes its effect on the learning-rate schedule and selected epoch. Table~\ref{tab:revision-loss} reports default greedy decoding for both conditions. This tests whether supervising the unused head helps the shared representation; it does not isolate each gradient pathway or redesign the loss around decoded cost.

% Generated by scripts/render_revision_tables.py.
\begin{table*}[!t]
\begin{revised}
\caption{\begin{revised}Matched log-loss ablation with default greedy decoding on 512 test rows. Full retains log supervision; no log sets only its coefficient to zero. All candidates pass replay. The difference is full minus no log in percentage points.\end{revised}}
\label{tab:revision-loss}
\centering
\renewcommand{\arraystretch}{1.15}
\begin{tabularx}{\linewidth}{@{}Yrrrrr@{}}
\toprule
Seed & \shortstack{Epoch\\full} & \shortstack{Epoch\\no log} & \shortstack{Optimal\\full} & \shortstack{Optimal\\no log} & \shortstack{Difference\\(points)}\\
\midrule
17 & 27 & 50 & 463 & 465 & -0.39\\
29 & 48 & 48 & 465 & 467 & -0.39\\
\bottomrule
\end{tabularx}
\end{revised}
\end{table*}


Removing log supervision yields 465 rather than 463 optima for seed 17 and 467 rather than 465 for seed 29. The full-minus-ablated difference is $-.39$ points for each seed, with paired family-bootstrap intervals $[-1.56,.78]$ and $[-1.17,.00]$. These small differences do not establish a benefit from log supervision for the default decoder. We retain the reference objective to preserve the original comparison, while reporting this negative ablation result and the separate use of log scores by repair.

Using the original seed-13 scores, single-log repair raises optimal candidates from 466 to 498, or from 91.02\% to 97.27\%, while mean gap falls from .180 to .039. Unguided repair improves from 445 to 479 optimal candidates, or from 86.91\% to 93.55\%, with mean gap .111. Thus additional candidate exploration helps substantially even without learning. Guided and unguided repair differ in both move rankings and trial order, so their difference cannot be attributed solely to the log head. The failure in Figure~\ref{fig:failure} is repaired from cost six to the optimum of two. The 6.25-point improvement over guided greedy decoding has a paired family-bootstrap interval of $[3.52,9.38]$ points, conditional on this checkpoint and test generator. The repaired method still misses 14 optima and remains a finite heuristic. Its additional computation is included in the timing comparison below.

\subsubsection{Structures Excluded from Training and a Second Compiler}

We generate 64 acyclic process trees and 64 trees with a loop at the root, using seed 703. Each class balances nominal visible sizes 6, 10, 14, and 18; sequence, choice, and parallel operators have weights .4, .3, and .3. The loop class adds a single visible redo leaf. An alphabet fraction of .7 permits repeated labels. Each tree supplies a clean playout and a corrupted observation at deviation rate .25; loop playout uses redo probability .5 and at most two redos. The same two observations are paired across the original block compiler and PM4Py's process-tree compiler, yielding 512 inputs.

Before accepting a tree, we exclude either compilation if its directed Petri-net graph is isomorphic to a training graph. The comparison preserves place/visible-transition/silent-transition types, arc weights, and initial/final tokens, while ignoring identifiers and activity strings. A structural hash filters candidates for full graph isomorphism checking. Two acyclic proposals are rejected. All training graphs are checked to be acyclic; both compilations of every loop input are checked to be cyclic. The loop class is therefore absent from training, and the PM4Py compiler provides an additional representation test. This does not exclude repeated visible languages or claim novelty relative to the validation split. Exact costs agree across all 256 paired observations; every clean observation has exact cost zero.

% Generated by scripts/render_revision_tables.py.
\begin{table*}[!t]
\begin{revised}
\caption{\begin{revised}Optimal candidate rates (\%) on structures excluded from training. Each row has 128 inputs from 64 trees and complete candidate/exact coverage. The three seeds use greedy decoding; guided repair uses seed 13.\end{revised}}
\label{tab:revision-structures}
\centering
\renewcommand{\arraystretch}{1.15}
\begin{tabularx}{\linewidth}{@{}Yrrrrrr@{}}
\toprule
Class, compiler & Unguided & Seed 13 & Seed 17 & Seed 29 & \shortstack{Repair\\unguided} & \shortstack{Repair\\guided}\\
\midrule
Acyclic, LARA & 76.6 & 81.2 & 81.2 & 78.1 & 82.0 & 86.7\\
Acyclic, PM4Py & 82.8 & 82.0 & 82.0 & 80.5 & 87.5 & 87.5\\
Loop, LARA & 54.7 & 52.3 & 51.6 & 52.3 & 64.1 & 65.6\\
Loop, PM4Py & 55.5 & 53.9 & 51.6 & 53.9 & 65.6 & 68.0\\
\bottomrule
\end{tabularx}
\end{revised}
\end{table*}


Table~\ref{tab:revision-structures} shows complete legal and exact coverage, but limited transfer of the learned advantage. On acyclic inputs, the three greedy checkpoints attain 78.1\% to 81.2\% optimality with the original compiler and 80.5\% to 82.0\% with PM4Py; unguided decoding attains 76.6\% and 82.8\%, respectively. On loops, all three checkpoints fall below their unguided control under both compilers. For seed 13, mean gap on loops is 4.156 with the original compiler and 4.797 with PM4Py; guided repair reduces these to 1.781 and 2.375. Repair improves candidates, but neither accepting new structures nor preserving their exact paired costs demonstrates a general learned advantage. The unfavorable cyclic result is evidence against that stronger claim.

\subsubsection{Held-Out Cases from Three Public Logs}

We add the Sepsis Cases event log~\cite{SepsisLog}, containing 1050 cases, 15,214 events, and 16 activity labels, to receipt and the existing road traffic sample. Only activity sequences enter the experiment. For each log, seed 37 fixes a 70\% discovery and 30\% evaluation split by case before discovery. The resulting discovery/evaluation counts are 70/30 for road traffic, 1003/431 for receipt, and 735/315 for sepsis. Inductive Miner uses only discovery cases at thresholds .0 and .5. All evaluation variants are retained: 6, 48, and 272 per threshold, respectively. Case sets are disjoint, but identical activity variants may occur in both sets. No target log is used for neural training or selection.

% Generated by scripts/render_revision_tables.py.
\begin{table*}[!t]
\begin{revised}
\caption{\begin{revised}Held-out real-log cases. Fractions report optimal counts over all evaluated variants or cases. Greedy columns use seed 13. Repair columns are variant counts; guided repair also uses seed 13. All 652 variant/model inputs have exact references and legal candidates.\end{revised}}
\label{tab:revision-real}
\centering
\renewcommand{\arraystretch}{1.15}
\begin{tabularx}{\linewidth}{@{}Yrrrrr@{}}
\toprule
Log & Threshold & \shortstack{Greedy\\variants} & \shortstack{Greedy\\cases} & \shortstack{Repair guided\\variants} & \shortstack{Repair unguided\\variants}\\
\midrule
Road traffic & 0.0 & 6/6 & 30/30 & 6/6 & 6/6\\
Road traffic & 0.5 & 6/6 & 30/30 & 6/6 & 6/6\\
Receipt & 0.0 & 40/48 & 423/431 & 40/48 & 40/48\\
Receipt & 0.5 & 33/48 & 398/431 & 39/48 & 39/48\\
Sepsis & 0.0 & 254/272 & 297/315 & 254/272 & 254/272\\
Sepsis & 0.5 & 263/272 & 297/315 & 267/272 & 269/272\\
\bottomrule
\end{tabularx}
\end{revised}
\end{table*}


Table~\ref{tab:revision-real} reports variant and case counts. Every greedy cost for all three training seeds equals the unguided cost on the same input. Guided repair improves receipt at threshold .5 from 33 to 39 optimal variants and from 398 to 421 optimal cases, exactly matching unguided repair. On sepsis at threshold .5, guided repair gives 267 optimal variants and 301 optimal cases, while unguided repair gives 269 variants and 303 cases. Other conditions have unchanged optimal counts under repair, although receipt at threshold .0 has a smaller mean gap. Thus the extra exploration helps some cases, but learning supplies no real-log quality advantage in this study.

The added experiment separates model discovery from evaluated cases and adds a third domain. It does not supply independent normative process models, a temporal split, or operational deployment evidence. The road traffic sample remains small. None of the six discovered nets has duplicate visible labels, so this test still lacks the real-world ambiguity with the clearest synthetic guidance benefit. Thresholds are repeated conditions within a log, not independent domains.

\subsubsection{Timing with Common Processing Boundaries}

After all training and quality jobs finish, a common harness reruns eight methods on the original 512 test rows. Each call starts with the same in-memory Petri net and trace and ends after parsing, when required, and final independent replay. Method-specific parameter construction, feature extraction, neural inference, candidate construction or search, and existing internal checks are included. Checkpoint loading, log loading, and model discovery are excluded for every method. Three warmup inputs per method precede five repetitions; input and method order are shuffled with seed 902. Every call has the same five-second wall-time budget. Approximation parameters match Section~\ref{sec:protocol}. The batch-based subset approximation is excluded because its unit of computation is a group of traces; the new timing comparison concerns per-trace methods only.

% Generated by scripts/render_revision_tables.py.
\begin{table*}[!t]
\begin{revised}
\caption{\begin{revised}Common-boundary CPU timings, including final independent replay. Each method has 512 inputs and five calls per input. Time summaries aggregate the per-input medians. Optimality uses all attempted calls as denominator.\end{revised}}
\label{tab:revision-timing}
\centering
\renewcommand{\arraystretch}{1.15}
\begin{tabularx}{\linewidth}{@{}Yrrrr@{}}
\toprule
Method & Legal calls & \shortstack{Optimal\\(\%)} & \shortstack{Median\\(ms)} & \shortstack{Interquartile range\\(ms)}\\
\midrule
Greedy, unguided & 2560/2560 & 86.91 & 0.307 & [0.275, 0.344]\\
Greedy, guided & 2560/2560 & 91.02 & 2.060 & [2.006, 2.118]\\
Repair, guided & 2560/2560 & 97.27 & 2.325 & [2.203, 2.467]\\
Exact A* & 2560/2560 & 100.00 & 0.664 & [0.579, 0.778]\\
Discounted A* & 2560/2560 & 95.70 & 0.414 & [0.364, 0.471]\\
Tandem repeats & 2560/2560 & 98.05 & 0.338 & [0.277, 0.417]\\
Sliding window & 2560/2560 & 100.00 & 0.330 & [0.269, 0.418]\\
Fixed horizon & 2560/2560 & 98.05 & 4.225 & [3.247, 5.248]\\
\bottomrule
\end{tabularx}
\end{revised}
\end{table*}


For each input and method, we first take the median of its five calls. Table~\ref{tab:revision-timing} reports the median and interquartile range of those 512 input medians, plus legal coverage and optimality over all 2560 attempted calls per method. The interquartile range describes differences across inputs, not a confidence interval. All 20,480 calls finish with legal outputs. Guided greedy decoding takes a median 2.060 ms, guided repair 2.325 ms, exact A* .664 ms, and sliding windows .330 ms. Unguided greedy decoding takes .307 ms. Thus the learned candidate is slower than exact search on these small inputs; repair adds a modest median cost while improving quality. Sliding windows are optimal throughout this test because its short traces fit inside one window. Fixed-horizon search is slower at 4.225 ms but has a higher optimal rate than guided repair. These measurements support a runtime comparison on this host and test distribution with equal outer boundaries. Candidate generation and exact search still provide different guarantees, and the experiment establishes no reduction in certified-mode computation or deployment latency.
\end{revised}
